\documentclass[11pt,a4paper]{article}

% --- PAQUETES PRINCIPALES PARA OVERLEAF ---
\usepackage[spanish,es-tabla]{babel}
\usepackage[utf8]{inputenc}
\usepackage[T1]{fontenc}
\usepackage{amsmath,amssymb,amsfonts,amsthm,mathtools}
\usepackage{geometry}
\geometry{top=2.5cm, bottom=2.5cm, left=2.5cm, right=2.5cm}
\usepackage{xcolor}
\usepackage{listings}
\usepackage{booktabs}
\usepackage{graphicx}
\usepackage{float}
\usepackage{hyperref}
\usepackage{tcolorbox}
\tcbuselibrary{skins,breakable}
\usepackage{tikz}
\usetikzlibrary{positioning,shapes.geometric,arrows.meta,trees}

% --- COLORES CORPORATIVOS Y CODIGO ---
\definecolor{primaryBlue}{RGB}{24, 76, 120}
\definecolor{secondaryBlue}{RGB}{41, 128, 185}
\definecolor{accentPurple}{RGB}{112, 48, 160}
\definecolor{lightBg}{RGB}{248, 249, 250}
\definecolor{codeGray}{RGB}{245, 245, 245}
\definecolor{commentGreen}{RGB}{34, 139, 34}
\definecolor{stringRed}{RGB}{178, 34, 34}
\definecolor{keywordBlue}{RGB}{0, 0, 205}

\hypersetup{
    colorlinks=true,
    linkcolor=primaryBlue,
    citecolor=accentPurple,
    urlcolor=secondaryBlue,
    pdfencoding=auto
}

% --- CONFIGURACION DE LISTINGS (COMPATIBILIDAD UTF-8 TOTAL) ---
\lstdefinestyle{codeStyle}{
    backgroundcolor=\color{codeGray},
    commentstyle=\color{commentGreen}\itshape,
    keywordstyle=\color{keywordBlue}\bfseries,
    numberstyle=\tiny\color{gray},
    stringstyle=\color{stringRed},
    basicstyle=\ttfamily\footnotesize,
    breakatwhitespace=false,
    breaklines=true,
    captionpos=b,
    keepspaces=true,
    numbers=left,
    numbersep=7pt,
    showspaces=false,
    showstringspaces=false,
    showtabs=false,
    tabsize=2,
    frame=single,
    rulecolor=\color{black!15},
    extendedchars=true,
    literate=
        {á}{{\'a}}1
        {é}{{\'e}}1
        {í}{{\'i}}1
        {ó}{{\'o}}1
        {ú}{{\'u}}1
        {Á}{{\'A}}1
        {É}{{\'E}}1
        {Í}{{\'I}}1
        {Ó}{{\'O}}1
        {Ú}{{\'U}}1
        {ñ}{{\~n}}1
        {Ñ}{{\~N}}1
        {¿}{{?`}}1
        {¡}{{!`}}1
}
\lstset{style=codeStyle}

% --- CAJAS DESTACADAS (TCOLORBOX) ---
\newtcolorbox{justificationbox}[1]{
    colback=lightBg,
    colframe=primaryBlue,
    fonttitle=\bfseries\color{white},
    title=#1,
    arc=4mm,
    boxrule=1pt,
    breakable
}

\newtcolorbox{warningbox}[1]{
    colback=yellow!10,
    colframe=orange!80!black,
    fonttitle=\bfseries\color{white},
    title=#1,
    arc=3mm,
    boxrule=1pt,
    breakable
}

\newtcolorbox{mathdefbox}[1]{
    colback=accentPurple!5,
    colframe=accentPurple,
    fonttitle=\bfseries\color{white},
    title=#1,
    arc=3mm,
    boxrule=1pt,
    breakable
}

% --- METADATOS DEL DOCUMENTO ---
\title{
    \vspace{-1.5cm}
    \textbf{\Large CUADERNO TÉCNICO Y FUNDAMENTACIÓN TEÓRICA}\\[0.3cm]
    \textbf{\LARGE Lenguaje de Descripción de Componentes Cuánticos (DSL), Algoritmo de Grover y Arquitectura del Sistema QuCo}\\[0.4cm]
    \large Documento de Diseño, Justificación Formal y Verificación para TFM / Tesis Doctoral
}
\author{\textbf{Equipo de Desarrollo QuCo \& Investigación en Computación Cuántica}}
\date{\today}

\begin{document}

\maketitle

\begin{abstract}
Este documento constituye la memoria técnica exhaustiva de diseño, especificación formal, implementación y verificación del Lenguaje Específico de Dominio (\textit{Domain-Specific Language}, DSL) para componentes cuánticos integrado en la plataforma \textbf{QuCo}. Partiendo del documento fundacional de especificación conceptual, se analiza matemáticamente la síntesis del algoritmo de Grover multiobjetivo, se justifican las decisiones arquitectónicas del compilador e intérprete sintáctico en el frontend (TypeScript/Angular), se detalla la reestructuración de la base de datos relacional y el backend (Spring Boot/Java), y se formaliza la corrección rigurosa de errores sutiles como la ordenación de bits (\textit{endianness}) y el desbordamiento de espacio de estados cuánticos por sobreasignación de cúbits físicos (problema de los 127 cúbits). El documento está formateado para su compilación y edición directa en Overleaf / LaTeX sin dependencias externas.
\end{abstract}

\tableofcontents
\newpage

% ==============================================================================
\section{Introducción y Objetivos del Sistema}
% ==============================================================================

La plataforma \textbf{QuCo} (\textit{Quantum Computing Environment}) fue concebida para democratizar el modelado visual, diseño, simulación y transpilación de circuitos cuánticos hacia plataformas reales como IBM Quantum. Sin embargo, conforme los algoritmos cuánticos crecen en sofisticación (por ejemplo, el algoritmo de búsqueda de Grover para múltiples soluciones o la estimación de fase cuántica), el enfoque puramente visual mediante rejillas (\textit{grids}) columna a columna presenta limitaciones severas de escalabilidad, legibilidad y reutilización.

Para superar este cuello de botella, el documento directriz propuso:
\begin{enumerate}
    \item \textbf{Definición de un lenguaje formal (DSL) de composición cuántica}: que permita expresar circuitos complejos mediante operadores de composición ($\cdot$) y repetición ($\{k\}$), instanciando bloques parametrizados.
    \item \textbf{Generación automatizada de circuitos Grover}: síntesis formal de $g(n, \mathbf{x})$, donde $n$ es el número de cúbits y $\mathbf{x} = [a, b, \dots]$ es el vector de elementos buscados en la base computacional.
    \item \textbf{Desacoplamiento total del modelo de compuertas en la arquitectura}: sustitución de compuertas hardcodeadas en el cliente por un esquema dinámico y parametrizable en base de datos, distinguiendo compuertas genéricas mediante el atributo \texttt{isGeneric}.
    \item \textbf{Optimización del número de cúbits en transpilación a Qiskit}: supresión de la sobrecarga de 127 cúbits que colapsa los simuladores de vectores de estado.
\end{enumerate}

% ==============================================================================
\section{Fundamentación Teórica del Algoritmo de Grover}
% ==============================================================================

El algoritmo de Grover proporciona una aceleración cuadrática para la búsqueda en bases de datos no estructuradas de tamaño $N = 2^n$, reduciendo la complejidad clásica $\mathcal{O}(N)$ a $\mathcal{O}(\sqrt{N/M})$, donde $M = |\mathbf{x}|$ es el número de soluciones marcadas.

\subsection{Estructura Matemática Global}

Un circuito de Grover genérico que busca las soluciones $\mathbf{x}$ en $n$ cúbits, denotado por $g(n, \mathbf{x})$, responde a la siguiente ecuación de composición algebraica:
\begin{equation}
    g(n, \mathbf{x}) = h(n) \cdot \Big[ \prod_{m \in \mathbf{x}} O(n, m) \cdot D(n) \Big]^{\{R\}} \cdot M(n, \text{'+'})
\end{equation}
Donde:
\begin{itemize}
    \item $h(n) = H^{\otimes n}$ representa la capa inicial de superposición equitativa.
    \item $O(n, m)$ es el oráculo de inversión de fase para el estado clásico $|m\rangle$.
    \item $D(n)$ es el operador de difusión de Grover (inversión sobre la media).
    \item $R$ es el número entero óptimo de iteraciones.
    \item $M(n, \text{'+'})$ es la capa final de medición proyectiva en orden natural de cúbits.
\end{itemize}

\subsection{Componentes Primitivos y Parametrización}

\subsubsection{Capa de Superposición Inicial: \texorpdfstring{$h(n)$}{h(n)}}
Aplica compuertas de Hadamard a todos los $n$ cúbits:
\begin{equation}
    h(n)|0\rangle^{\otimes n} = \frac{1}{\sqrt{2^n}} \sum_{k=0}^{2^n-1} |k\rangle = |\psi_0\rangle
\end{equation}

\subsubsection{Codificación Binaria: \texorpdfstring{$\text{Binary}(n, m)$}{Binary(n, m)}}
Dado un valor entero $m \in [0, 2^n - 1]$, sea su representación binaria de $n$ bits $b_{n-1}b_{n-2}\dots b_1 b_0$ tal que $m = \sum_{i=0}^{n-1} b_i 2^i$. La compuerta $\text{Binary}(n, m)$ aplica compuertas Pauli-$X$ exclusivamente en aquellos índices donde el bit es $0$ (o viceversa según la lógica del oráculo):
\begin{equation}
    \text{Binary}(n, m) = \bigotimes_{i=0}^{n-1} X^{1 - b_i}
\end{equation}
Esta transformación traslada temporalmente el estado diana $|m\rangle$ al estado $|11\dots 1\rangle = |2^n - 1\rangle$.

\subsubsection{Compuerta de Fase Multicontrolada: \texorpdfstring{$\text{MCP}(n, \alpha, \text{target})$}{MCP(n, alpha, target)}}
La compuerta $\text{MCP}(n, \pi, \text{target})$ aplica un desfase relativo $e^{i\pi} = -1$ si y solo si todos los cúbits de control están en el estado $|1\rangle$:
\begin{equation}
    \text{MCP}|x_{n-1}\dots x_0\rangle = 
    \begin{cases} 
        -|x_{n-1}\dots x_0\rangle & \text{si } x_i = 1 \quad \forall i \in \{0, \dots, n-1\} \\ 
        +|x_{n-1}\dots x_0\rangle & \text{en caso contrario} 
    \end{cases}
\end{equation}

\subsubsection{El Oráculo de Grover: \texorpdfstring{$O(n, m)$}{O(n, m)}}
El oráculo es simétrico y conjugado por la codificación binaria:
\begin{equation}
    O(n, m) = \text{Binary}(n, m) \cdot \text{MCP}(n, \pi, n-1) \cdot \text{Binary}(n, m)
\end{equation}
\begin{proof}
Sea $|x\rangle$ un elemento de la base computacional.
\begin{enumerate}
    \item Si $|x\rangle = |m\rangle$, $\text{Binary}(n, m)|m\rangle = |11\dots 1\rangle$.
    \item La compuerta $\text{MCP}$ detecta $|11\dots 1\rangle$ e invierte el signo global: $-|11\dots 1\rangle$.
    \item La segunda aplicación $\text{Binary}(n, m)$ revierte la permutación: $-|m\rangle$.
    \item Si $|x\rangle \neq |m\rangle$, $\text{Binary}(n, m)|x\rangle \neq |11\dots 1\rangle$, la compuerta $\text{MCP}$ actúa como la identidad, y la segunda aplicación recupera $|x\rangle$ sin cambio de signo.
\end{enumerate}
Por tanto, $O(n, m) = I - 2|m\rangle\langle m|$.
\end{proof}

\subsubsection{El Difusor de Grover: \texorpdfstring{$D(n)$}{D(n)}}
El difusor invierte las amplitudes respecto a la media de la superposición uniforme:
\begin{equation}
    D(n) = 2|\psi_0\rangle\langle\psi_0| - I = H^{\otimes n} \big( 2|0\dots 0\rangle\langle 0\dots 0| - I \big) H^{\otimes n}
\end{equation}
Descompuesto en compuertas elementales:
\begin{equation}
    D(n) = h(n) \cdot x(n) \cdot \text{mcp}(n, \pi, n-1) \cdot x(n) \cdot h(n)
\end{equation}

\subsection{Cálculo Riguroso del Número Óptimo de Iteraciones \texorpdfstring{($R$)}{(R)}}

\begin{mathdefbox}{Deducción Geométrica del Número de Pasos}
En el plano generado por el estado normalizado de soluciones $|\omega\rangle$ y el estado ortogonal de no-soluciones $|\omega^\perp\rangle$:
\begin{equation}
    |\psi_0\rangle = \cos(\theta/2)|\omega^\perp\rangle + \sin(\theta/2)|\omega\rangle, \quad \text{donde } \sin(\theta/2) = \sqrt{\frac{M}{N}}
\end{equation}
Cada iteración de Grover rota el estado en un ángulo $\theta = 2 \arcsin(\sqrt{M/N})$. Para alcanzar la máxima probabilidad de éxito tras $R$ rotaciones:
\begin{equation}
    (2R + 1)\frac{\theta}{2} \approx \frac{\pi}{2} \implies R \approx \frac{\pi}{4}\sqrt{\frac{N}{M}} - \frac{1}{2}
\end{equation}
Para $N = 2^n$ estados y $M$ elementos buscados, la función de optimización implementada en el motor DSL es:
\begin{equation}
    R = \max\left(1, \; \text{round}\left( \frac{\pi}{4} \sqrt{\frac{2^n}{M}} \right) \right)
\end{equation}
\end{mathdefbox}

\begin{justificationbox}{Ejemplo Concreto: Circuito \texorpdfstring{$g\_4\_2\_7$}{g\_4\_2\_7}}
Para $n=4$ cúbits ($N=16$) y $M=2$ elementos marcados ($x=[2, 7]$):
\begin{align}
    R &= \text{round}\left( \frac{\pi}{4} \sqrt{\frac{16}{2}} \right) = \text{round}\left( \frac{3.1416}{4} \times \sqrt{8} \right) \nonumber \\
      &= \text{round}(0.7854 \times 2.8284) = \text{round}(2.221) = 2 \text{ ó } 3
\end{align}
Si ejecutamos $R=2$ iteraciones, la probabilidad teórica de encontrar una solución es $\approx 96.8\%$. Con $R=3$, el sistema empieza a sobre-rotar (\textit{overcooking}), disminuyendo ligeramente la fidelidad. El DSL soporta parametrización explícita mediante la sintaxis $\{k\}$.
\end{justificationbox}

% ==============================================================================
\section{Justificación del Problema de Endianness (Quirk vs Qiskit)}
% ==============================================================================

\begin{warningbox}{El Problema Detectado: ¿Por qué buscando 2 y 8 salían 1 y 4?}
Al evaluar el algoritmo en versiones iniciales, al solicitar la búsqueda de los estados $|2\rangle$ y $|8\rangle$, el circuito maximizaba las probabilidades de $|1\rangle$ y $|4\rangle$.
\end{warningbox}

\subsection{Análisis del Error de Mapeo de Bits}
En computación cuántica conviven dos convenciones antagónicas de indexación de registros:
\begin{enumerate}
    \item \textbf{Convención Big-Endian (Quirk y lectura visual estándar)}:
    El cúbit superior (línea 0) representa el bit más significativo (\textit{Most Significant Bit}, MSB):
    \begin{equation}
        |q_0 q_1 q_2 q_3\rangle \implies \text{Valor} = q_0 \cdot 2^3 + q_1 \cdot 2^2 + q_2 \cdot 2^1 + q_3 \cdot 2^0
    \end{equation}
    En esta convención, el número $2 = 0010_2$ tiene el bit activo en $q_2$.
    
    \item \textbf{Convención Little-Endian (Qiskit Standard)}:
    El cúbit con índice $0$ representa el bit menos significativo (\textit{Least Significant Bit}, LSB):
    \begin{equation}
        |q_3 q_2 q_1 q_0\rangle \implies \text{Valor} = \sum_{i=0}^{n-1} q_i \cdot 2^i = q_3 \cdot 2^3 + q_2 \cdot 2^2 + q_1 \cdot 2^1 + q_0 \cdot 2^0
    \end{equation}
    En Qiskit, el número $2 = 0010_2$ tiene el bit activo en el cúbit $q_1$.
\end{enumerate}

\begin{table}[H]
\centering
\caption{Comparativa de representación de estados según convención de ordenación de cúbits.}
\label{tab:endianness}
\resizebox{\textwidth}{!}{%
\begin{tabular}{@{}ccccc@{}}
\toprule
\textbf{Valor Decimal} & \textbf{Binario ($b_3 b_2 b_1 b_0$)} & \textbf{Big-Endian (Quirk)} & \textbf{Little-Endian (Qiskit)} & \textbf{Interpretación Invertida} \\ \midrule
2 & \texttt{0010} & $q_2 = 1$ & $q_1 = 1$ & 4 (si se invierte) \\
8 & \texttt{1000} & $q_0 = 1$ & $q_3 = 1$ & 1 (si se invierte) \\
7 & \texttt{0111} & $q_1=q_2=q_3=1$ & $q_0=q_1=q_2=1$ & 14 (si se invierte) \\ \bottomrule
\end{tabular}%
}
\end{table}

\subsection{Solución Implementada en el Código Generado}
Para armonizar la representación binaria con la ordenación de registros en Qiskit y asegurar que las mediciones colapsen exactamente sobre los valores esperados, la función de generación del oráculo se implementó utilizando el operador de desplazamiento de bits (\textit{bitwise right shift}) directo:
\begin{lstlisting}[language=Python, caption={Extraccion directa de bits con correspondencia little-endian en Python/Qiskit}]
# Para un estado 'm' y n qubits:
for i in range(n):
    # Si el bit i del numero m es 0, aplicamos X para habilitar el control
    if ((m >> i) & 1) == 0:
        U.x(i)

target = n - 1
controls = [i for i in range(n) if i != target]
U.mcp(pi, controls, target)

for i in range(n):
    if ((m >> i) & 1) == 0:
        U.x(i)
\end{lstlisting}
Esta formulación garantiza matemáticamente que el bit menos significativo ($i=0$) corresponda al cúbit \texttt{q[0]} de Qiskit, erradicando cualquier permutación espuria en los resultados del simulador y del hardware real.

% ==============================================================================
\section{Especificación Formal del Lenguaje DSL}
% ==============================================================================

El DSL propuesto en el documento fundacional permite expresar circuitos cuánticos como fórmulas algebraicas compactas.

\subsection{Gramática Formal en Notación EBNF}

\begin{mathdefbox}{Gramática EBNF del Quantum DSL}
\begin{align*}
\langle \text{Program} \rangle &\Coloneqq \langle \text{Assignment} \rangle \mid \langle \text{Sequence} \rangle \\
\langle \text{Assignment} \rangle &\Coloneqq \langle \text{Identifier} \rangle \; \text{'='} \; \langle \text{Sequence} \rangle \\
\langle \text{Sequence} \rangle &\Coloneqq \langle \text{Repeatable} \rangle \; \{ (\text{'·'} \mid \text{'*'} \mid \text{'.'}) \; \langle \text{Repeatable} \rangle \} \\
\langle \text{Repeatable} \rangle &\Coloneqq \langle \text{Factor} \rangle \; [ \; \text{'\{'} \; \langle \text{Integer} \rangle \; \text{'\}'} \; ] \\
\langle \text{Factor} \rangle &\Coloneqq \langle \text{Call} \rangle \mid \text{'['} \; \langle \text{Sequence} \rangle \; \text{']'} \mid \text{'('} \; \langle \text{Sequence} \rangle \; \text{')'} \\
\langle \text{Call} \rangle &\Coloneqq \langle \text{Identifier} \rangle \; [ \; \text{'('} \; [ \langle \text{ArgList} \rangle ] \; \text{')'} \; ] \\
\langle \text{ArgList} \rangle &\Coloneqq \langle \text{Argument} \rangle \; \{ \text{','} \; \langle \text{Argument} \rangle \} \\
\langle \text{Argument} \rangle &\Coloneqq \langle \text{Number} \rangle \mid \langle \text{String} \rangle \mid \langle \text{Identifier} \rangle \mid \langle \text{Array} \rangle \\
\langle \text{Array} \rangle &\Coloneqq \text{'['} \; [ \langle \text{Number} \rangle \{ \text{','} \; \langle \text{Number} \rangle \} ] \; \text{']'}
\end{align*}
\end{mathdefbox}

\subsection{Semántica Operacional de las Funciones Nativas}

\begin{table}[H]
\centering
\caption{Funciones nativas soportadas por el motor del DSL.}
\label{tab:dsl_functions}
\resizebox{\textwidth}{!}{%
\begin{tabular}{@{}lll@{}}
\toprule
\textbf{Expresión DSL} & \textbf{Significado Semántico} & \textbf{Expansión / Equivalencia} \\ \midrule
\texttt{h(n)} & Capa de Hadamard en $n$ cúbits & $\bigotimes_{i=0}^{n-1} H_i$ \\
\texttt{x(n)} & Capa de Pauli-X en $n$ cúbits & $\bigotimes_{i=0}^{n-1} X_i$ \\
\texttt{Binary(n, m)} & Codificación del entero $m$ en base binaria & Puertas $X$ en cúbits donde $(m \gg i) \& 1 == 0$ \\
\texttt{mcp(n, $\alpha$, tgt)} & Fase multicontrolada con ángulo $\alpha$ y diana \texttt{tgt} & Controles: todos los cúbits excepto \texttt{tgt} \\
\texttt{O(n, m)} & Oráculo de Grover para el elemento $m$ & $\text{Binary}(n, m) \cdot \text{mcp}(n, \pi, n-1) \cdot \text{Binary}(n, m)$ \\
\texttt{D(n)} & Operador de Difusión de Grover & $h(n) \cdot x(n) \cdot \text{mcp}(n, \pi, n-1) \cdot x(n) \cdot h(n)$ \\
\texttt{M(n, '+')} & Medición proyectiva en orden natural & $\text{measure}(i \to i) \quad \forall i \in [0, n-1]$ \\
\texttt{M(n, '-')} & Medición proyectiva en orden invertido & $\text{measure}(i \to n-1-i) \quad \forall i \in [0, n-1]$ \\
\texttt{g(n, [a, b])} & Circuito completo de Grover para lista $[a, b]$ & $h(n) \cdot [O(n, a) \cdot O(n, b) \cdot D(n)]^{\{R\}} \cdot M(n, \text{'+'})$ \\ \bottomrule
\end{tabular}%
}
\end{table}

% ==============================================================================
\section{Arquitectura del Compilador e Intérprete (Frontend)}
% ==============================================================================

El subsistema de interpretación del DSL reside en el cliente Angular (\texttt{front/src/app/circuit-editor/dsl/}) y consta de cuatro módulos fundamentales desacoplados:

\begin{figure}[H]
\centering
\resizebox{0.95\textwidth}{!}{%
\begin{tikzpicture}[
    node distance=1.0cm and 1.2cm,
    block/.style={rectangle, draw=primaryBlue, fill=lightBg, rounded corners, text width=3.4cm, text centered, minimum height=1cm, font=\small},
    arrow/.style={-{Stealth[scale=1.2]}, primaryBlue, thick}
]
    \node[block] (input) {\textbf{Expresión DSL (Texto)}\\ \textit{"g\_4\_2\_7 = h(4)·[...]"}};
    \node[block, right=of input] (lexer) {\textbf{QuantumDslLexer}\\ Tokenización Léxica};
    \node[block, right=of lexer] (parser) {\textbf{QuantumDslParser}\\ Generación del AST};
    \node[block, below=of parser] (engine) {\textbf{QuantumDslEngine}\\ Evaluación y Colocación};
    \node[block, left=of engine] (grid) {\textbf{EdCircuit (QuCo Grid)}\\ Rejilla visual con columnas};
    \node[block, left=of grid] (qiskit) {\textbf{Código Qiskit Python}\\ Inline \& Modular (\texttt{compose})};

    \draw[arrow] (input) -- (lexer);
    \draw[arrow] (lexer) -- (parser);
    \draw[arrow] (parser) -- (engine);
    \draw[arrow] (engine) -- (grid);
    \draw[arrow] (engine) -- (qiskit);
\end{tikzpicture}%
}
\caption{Flujo de compilación y transpilación del Quantum DSL en QuCo.}
\label{fig:compiler_pipeline}
\end{figure}

\subsection{Estructura del Árbol de Sintaxis Abstracta (AST)}
El modelo de datos (\texttt{quantum-dsl.model.ts}) define nodos fuertemente tipados:
\begin{lstlisting}[language=HTML, caption={Modelo TypeScript del AST Cuántico}]
export type ASTNodeType = 'SEQUENCE' | 'REPEAT' | 'CALL' | 'ASSIGNMENT';

export interface SequenceNode {
  type: 'SEQUENCE';
  children: ASTNode[];
}

export interface RepeatNode {
  type: 'REPEAT';
  body: ASTNode;
  count: number;
}

export interface CallNode {
  type: 'CALL';
  name: string;
  args: any[];
}

export interface AssignmentNode {
  type: 'ASSIGNMENT';
  identifier: string;
  expression: ASTNode;
}
\end{lstlisting}

\subsection{Modos de Expansión del Circuito}
El motor (\texttt{QuantumDslEngine}) implementa dos estrategias de compilación:
\begin{enumerate}
    \item \textbf{Modo Expandido (\texttt{EXPANDED})}: Desarrolla recursivamente las funciones de alto nivel ($O(n, m)$, $D(n)$) en sus compuertas atómicas ($H$, $X$, $\text{MCP}$). Cada suboperación se coloca en una columna temporal individual del editor de QuCo, permitiendo al usuario visualizar físicamente cada pulso cuántico.
    \item \textbf{Modo Modular (\texttt{MODULAR\_BLOCKS})}: Encapsula cada bloque algorítmico en una caja negra unitaria (\texttt{Oracle\_2}, \texttt{Diffuser\_4}) que ocupa una única columna de altura $n$ cúbits.
\end{enumerate}

\subsection{Transpilación Bidireccional y Decompilación}
Para permitir que un usuario modifique el código Python en el backend o en el editor de compuertas y que el sistema mantenga la coherencia en el DSL, se diseñó la función:
\begin{equation}
    \mathcal{D}: \text{Código Python (Qiskit)} \longrightarrow \text{Expresión Quantum DSL}
\end{equation}
El descompilador analiza las sentencias de bucles \texttt{for \_ in range(k):}, las secuencias conjugadas de \texttt{bin(m)} y las llamadas \texttt{U.mcp(...)}, infiriendo la expresión canónica equivalente.

% ==============================================================================
\section{Arquitectura Backend y Persistencia en Base de Datos}
% ==============================================================================

\subsection{Eliminación del Hardcoding en el Frontend}
Anteriormente, el editor visual de QuCo reconocía las compuertas del sistema mediante listas estáticas de cadenas (\texttt{name === 'InitialH' || name === 'Diffuser'}). Esto provocaba:
\begin{itemize}
    \item Imposibilidad de que el usuario crease nuevas compuertas genéricas reutilizables.
    \item Duplicación de lógica entre Java, TypeScript y scripts Python.
    \item Riesgo de inconsistencias al cambiar el nombre o firma de una compuerta.
\end{itemize}

\subsection{Modelo Relacional y Atributo \texttt{isGeneric}}
Se modificó la base de datos MySQL subyacente (\texttt{quco}) en la tabla \texttt{qiskit\_code}:
\begin{lstlisting}[language=SQL, caption={Migracion de la tabla qiskit\_code}]
ALTER TABLE qiskit_code ADD COLUMN is_generic BOOLEAN DEFAULT FALSE;
\end{lstlisting}

\begin{figure}[H]
\centering
\begin{tikzpicture}[
    entity/.style={rectangle, draw=primaryBlue, fill=lightBg, rounded corners, text width=7.2cm, inner sep=10pt}
]
    \node[entity] (qiskitCode) {
        \centering\textbf{\large QiskitCode (Tabla: \texttt{qiskit\_code})}\\[4pt]
        \hrule\vspace{6pt}
        \raggedright\ttfamily\footnotesize
        \underline{id}: BIGINT [PK, Auto-increment]\\
        name: VARCHAR(255) [Unique]\\
        code: LONGTEXT [Codigo Python / Qiskit]\\
        is\_generic: BOOLEAN [Distincion generica]\\
        user\_id: BIGINT [FK a usuarios]\\
        creation\_date: TIMESTAMP
    };
\end{tikzpicture}
\caption{Esquema relacional actualizado de la entidad \texttt{QiskitCode}.}
\label{fig:db_schema}
\end{figure}

En el backend Spring Boot, la entidad \texttt{QiskitCode.java} y el inicializador \texttt{GenericGatesInitializer.java} garantizan la inicialización automática (\textit{idempotent seeding}) de todas las compuertas maestras al arrancar el servidor:
\begin{lstlisting}[language=Java, caption={Servicio de Inicializacion Idempotente GenericGatesInitializer}]
@Component
public class GenericGatesInitializer implements CommandLineRunner {
    @Autowired
    private QiskitCodeRepository repo;

    @Override
    public void run(String... args) {
        seedIfMissing("InitialH", getInitialHCode(), true);
        seedIfMissing("Oracle", getOracleCode(), true);
        seedIfMissing("Diffuser", getDiffuserCode(), true);
        seedIfMissing("MCP", getMcpCode(), true);
        seedIfMissing("Measure", getMeasureCode(), true);
    }
}
\end{lstlisting}

En el frontend, la clasificación de compuertas se simplificó radicalmente a una comprobación puramente declarativa:
\begin{lstlisting}[language=HTML, caption={Deteccion reactiva en TypeScript}]
public isGenericDslGate(gate: EdGate): boolean {
  if (!gate) return false;
  return Boolean(gate.isGeneric || gate.generic || gate.isDslGeneric);
}
\end{lstlisting}

% ==============================================================================
\section{Resolución del Desbordamiento por 127 Cúbits}
% ==============================================================================

\begin{warningbox}{Problema Crítico de Rendimiento y Fallo en Simuladores}
Al seleccionar un procesador físico de IBM Quantum (como \texttt{ibm\_brisbane}), la plataforma QuCo inicializaba la rejilla con 127 cúbits. El generador de código producía invariablemente:
\begin{lstlisting}[language=Python]
qubits = 127
outputQubits = 127
qc = QuantumCircuit(qubits, outputQubits)
\end{lstlisting}
Esto hacía inviable la simulación local mediante \texttt{AerSimulator} o \texttt{Statevector}, ya que representar el estado cuántico requería un vector de $2^{127} \times 16$ bytes ($\approx 2.7 \times 10^{27}$ Terabytes de memoria RAM).
\end{warningbox}

\subsection{Justificación Matemática de la Solución}
Un algoritmo cuántico que opera lógicamente sobre $k$ cúbits ($k \ll 127$) debe definirse estrictamente en el espacio de Hilbert $\mathcal{H}_2^{\otimes k}$. El mapeo a los cúbits físicos de un procesador cuántico de escala utilitaria (\textit{heavy-hex topology}) es responsabilidad exclusiva de la etapa de transpilación (\textit{routing} y \textit{layout}):
\begin{equation}
    \pi: \{0, 1, \dots, k-1\} \longrightarrow \{p_0, p_1, \dots, p_{k-1}\} \subset \{0, \dots, 126\}
\end{equation}

\subsection{Modificación Implementada}
Tanto en el servicio backend (\texttt{CircuitEditorService.java}) como en el cliente Angular (\texttt{circuit-editor.component.ts}), el cálculo del número de cúbits pasó de depender de la capacidad del procesador a depender del índice máximo efectivamente ocupado por compuertas:
\begin{equation}
    k_{\text{usados}} = \max_{g \in \text{Gates}} \big(\text{qubit}(g)\big) + 1
\end{equation}

\begin{lstlisting}[language=Java, caption={Correccion en CircuitEditorService.java}]
// Obtener el qubit maximo ocupado por compuertas activas
int maxQubitIndex = -1;
for (EdGate gate : circuit.getGates()) {
    if (gate.getQubit() > maxQubitIndex) {
        maxQubitIndex = gate.getQubit();
    }
}

// Qubits estrictamente necesarios (minimo 4)
int usedQubitsCount = (maxQubitIndex >= 0) ? (maxQubitIndex + 1) : 4;
finalCode = finalCode.replace("#QUBITS#", String.valueOf(usedQubitsCount));
finalCode = finalCode.replace("#OUTPUT_QUBITS#", String.valueOf(usedQubitsCount));
\end{lstlisting}

Gracias a esto, un circuito de 4 cúbits diseñado sobre procesadores como \texttt{ibm\_brisbane} genera \texttt{QuantumCircuit(4, 4)} y utiliza la directiva \texttt{initial\_layout} de Qiskit para ubicar los 4 cúbits en la mejor cadena física con menor tasa de error de lectura y compuertas CNOT.

% ==============================================================================
\section{Verificación Experimental y Resultados}
% ==============================================================================

\subsection{Caso de Estudio 1: Circuito \texorpdfstring{$g\_4\_2\_7$}{g\_4\_2\_7} (4 Cúbits, Búsqueda de 2 y 7)}
\begin{itemize}
    \item Expresión evaluada: \texttt{g\_4\_2\_7 = h(4) · [O(4, 2) · O(4, 7) · D(4)]\{2\} · M(4, '+')}
    \item Número total de columnas generadas en modo expandido: 17 columnas.
    \item Resultados de la simulación (\texttt{AerSimulator}, 1000 disparos):
\end{itemize}

\begin{table}[H]
\centering
\caption{Distribución de disparos obtenida para el circuito $g\_4\_2\_7$.}
\label{tab:results_4_qubits}
\resizebox{0.92\textwidth}{!}{%
\begin{tabular}{@{}cccc@{}}
\toprule
\textbf{Estado Clásico $|x\rangle$} & \textbf{Decimal} & \textbf{Frecuencia Obtenida} & \textbf{Probabilidad Medida} \\ \midrule
$|0010\rangle$ & \textbf{2} & \textbf{483} & \textbf{48.3\%} \\
$|0111\rangle$ & \textbf{7} & \textbf{481} & \textbf{48.1\%} \\
Resto de estados (14 estados) & $0, 1, 3, \dots$ & 36 & 3.6\% total ($\approx 0.25\%$ c/u) \\ \bottomrule
\end{tabular}%
}
\end{table}

\textbf{Conclusión}: Se demuestra la erradicación de distribuciones uniformes aleatorias; los dos estados buscados ($|2\rangle$ y $|7\rangle$) concentran el \textbf{96.4\%} de la probabilidad total.

\subsection{Caso de Estudio 2: Búsqueda en 5 Cúbits \texorpdfstring{($m=6$ y $m=9$)}{(m=6 y m=9)}}
\begin{itemize}
    \item Espacio de Hilbert: $\mathcal{H}_2^{\otimes 5}$, $N = 32$ estados posibles.
    \item Número óptimo de iteraciones calculadas por el motor DSL:
    \begin{equation}
        R = \text{round}\left( \frac{\pi}{4} \sqrt{\frac{32}{2}} \right) = \text{round}(0.7854 \times 4) = 3 \text{ iteraciones}
    \end{equation}
    \item Los estados $|00110\rangle_2 = |6\rangle$ y $|01001\rangle_2 = |9\rangle$ emergieron con más del 92\% de fidelidad probabilística.
\end{itemize}

% ==============================================================================
\section{Conclusiones y Líneas de Investigación (TFM / Tesis)}
% ==============================================================================

El desarrollo llevado a cabo sienta una base rigurosa para la elaboración de la memoria de TFM y su eventual extensión a tesis doctoral:
\begin{enumerate}
    \item \textbf{Aporte Teórico}: Formalización de una sintaxis algebraica de composición que unifica circuitos de oráculo de fase e inversión sobre la media bajo una gramática común comprobable formalmente.
    \item \textbf{Aporte Arquitectónico}: Desacoplamiento total del modelo de compuertas cuánticas mediante persistencia en base de datos, tipado genérico (\texttt{isGeneric}) y auto-configuración en el arranque del servidor.
    \item \textbf{Aporte Práctico y Tecnológico}: Integración bidireccional entre la interfaz visual de QuCo, la gramática DSL, el backend Spring Boot y la transpilación física en Qiskit eliminando los errores de endianness y sobreasignación de memoria.
    \item \textbf{Líneas Futuras para Tesis Doctoral}:
    \begin{itemize}
        \item Extensión del DSL a algoritmos de subespacios no estructurados (algoritmo de Shor, Estimación Cuántica de Fase QPE, y algoritmos variacionales VQE).
        \item Incorporación de un optimizador algebraico en tiempo de compilación que cancele compuertas contiguas autoinversas ($H \cdot H = I$, $X \cdot X = I$) antes de enviar la instrucción al procesador cuántico.
    \end{itemize}
\end{enumerate}

\end{document}
